%! Complex Numbers and the Library of Parent Functions — Unit 2, Lesson 2.1 — Homework (ANSWER KEY)
% THE LESSON'S GRADED PRACTICE and the last component in the packet — exactly TWO pages, started
% in the Homework Launch (the last 3 minutes of the period) and finished at home. Every context
% here is new: the notes used 7 - 2i, (2 + 3i)(4 - i), a ten-panel gallery, |x - 2| and
% |x + 3|; the AP Practice used (2 - i)(3 + 4i) and a machined rod; the homework uses new
% numbers, four new unlabeled graphs, |x - 4| and |x + 1|.
%   Page 1 — Part A (imaginary and complex numbers, 1–4) and Part B (the library, 5).
%   Page 2 — Part B continued (6–7), Part C (absolute value as a piecewise function, 8–9) and Part D (the headline, 10).
% The diagnostic item is 8: |x - 4| at x = 1 — does the student put -(x - 4) on the side
% where x - 4 is negative, and get +3? Problem 10 asks for the lesson's headline in writing.
%
% PAGE LOCKSTEP with the other half of the pair: work blocks are byte-identical, prose answers
% are \writespace{H}{…} (fixed height both ways), and each \blank sits at the END of a line or
% before a \qquad so its \ans cannot rewrap anything. The explicit \newpage fixes the page
% break in both files.
\documentclass[10pt]{article}
\usepackage{precalculus-article}
\usepackage{precalculus-key}
\newcommand{\TallMath}[1]{$\displaystyle #1\rule[-1.4em]{0pt}{3.2em}$}
% One small unlabeled graph for problem 5: #1 letter, #2 the plot (clipped to the window).
\newcommand{\hwgraph}[2]{%
  \begin{minipage}[t]{0.23\linewidth}\centering
  {\small\bfseries\color{plum}#1}\par\vspace{2pt}
  \begin{tikzpicture}[x=0.3cm, y=0.3cm]
    \draw[linegray, very thin, step=1] (-4,-4) grid (4,4);
    \draw[->, thin] (-4.3,0) -- (4.5,0);
    \draw[->, thin] (0,-4.3) -- (0,4.5);
    \begin{scope}
      \clip (-4,-4) rectangle (4,4);
      #2
    \end{scope}
  \end{tikzpicture}
  \end{minipage}}

\begin{document}

\pageheader{Unit 2, Lesson 2.1}{Homework --- Answer Key}

\vspace{0.12in}

\boxguard[8]
\begin{remindbox}
\small
\textbf{This is your graded homework.} Start it in class; finish it on your own by the due
date on your cover sheet. Show your work: rewrite every square root of a negative with $i$
\emph{before} you compute, and on any absolute value decide \emph{which rule owns the input}
first.
\end{remindbox}

\vspace{0.12in}

\begin{headlinebox}{lilac}
\textbf{\color{plum}Part A --- Imaginary and complex numbers}
\end{headlinebox}

\begin{enumerate}[leftmargin=1.9em, itemsep=7pt]

\item Write each using $i$, and simplify.

      $\sqrt{-49} = $ \ans{$7i$} \qquad $\sqrt{-20} = $ \ans{$2i\sqrt{5}$} \qquad
      $i^{9} = $ \ans{$i$}

\item Solve each equation.

      \begin{minipage}[t]{0.47\linewidth}
      \textbf{(a)} $x^2 + 64 = 0$
\begin{work}
  x^2 &= -64 \\
    x &= \pm\sqrt{-64} = \pm 8i
\end{work}
      \end{minipage}\hfill
      \begin{minipage}[t]{0.47\linewidth}
      \textbf{(b)} $x^2 = -50$
\begin{work}
  x &= \pm\sqrt{-50} \\
    &= \pm i\sqrt{50} = \pm 5i\sqrt{2}
\end{work}
      \end{minipage}

\item Simplify. Write each answer in the form $a + bi$.

      \begin{minipage}[t]{0.47\linewidth}
      \textbf{(a)} $(6 - 2i) - (1 + 5i)$
\begin{work}
  (6 - 2i) - (1 + 5i) &= 6 - 2i - 1 - 5i \\
                      &= 5 - 7i
\end{work}
      \end{minipage}\hfill
      \begin{minipage}[t]{0.47\linewidth}
      \textbf{(b)} $(5 + 2i)(3 - i)$
\begin{work}
  (5 + 2i)(3 - i) &= 15 - 5i + 6i - 2i^2 \\
                  &= 15 + i + 2 \\
                  &= 17 + i
\end{work}
      \end{minipage}

\item Write the complex conjugate of $1 + 3i$, then multiply the pair. What kind of number is
      the product, and why did that happen?
      \writespace{1.3cm}{$1 - 3i$. $(1 + 3i)(1 - 3i) = 1 - 9i^2 = 1 + 9 = 10$ --- a real number: the $3i$ and $-3i$ terms cancel, and $i^2 = -1$.}

\end{enumerate}

\vspace{0.06in}

\begin{headlinebox}{lilac}
\textbf{\color{plum}Part B --- The library of parent functions}
\end{headlinebox}

\begin{enumerate}[leftmargin=1.9em, itemsep=7pt, start=5]

\item Each graph below is a parent function. Write the letter of the matching graph.

      \smallskip
      \hwgraph{A}{%
        \draw[very thick, plum] plot[domain=0:4, samples=60, smooth] (\x,{sqrt(\x)});}\hfill
      \hwgraph{B}{%
        \draw[very thick, plum] plot[domain=-1.62:1.62, samples=50, smooth] (\x,{\x*\x*\x});}\hfill
      \hwgraph{C}{%
        \draw[very thick, plum] plot[domain=0.25:4, samples=50, smooth] (\x,{1/\x});
        \draw[very thick, plum] plot[domain=-4:-0.25, samples=50, smooth] (\x,{1/\x});}\hfill
      \hwgraph{D}{%
        \draw[very thick, plum] (-4,4) -- (0,0) -- (4,4);}

      \smallskip
      $y = |x|$: \ans{D} \qquad $y = x^3$: \ans{B} \qquad
      $y = \dfrac{1}{x}$: \ans{C} \qquad $y = \sqrt{x}$: \ans{A}

\end{enumerate}

\newpage

\begin{enumerate}[leftmargin=1.9em, itemsep=7pt, start=6]

\item \begin{minipage}[t]{0.48\linewidth}
      \textbf{(a)} Name every parent function in the library whose range is $y \ge 0$.
      \writespace{1.0cm}{$y = x^2$, $y = x^4$, $y = |x|$, $y = \sqrt{x}$}
      \end{minipage}\hfill
      \begin{minipage}[t]{0.48\linewidth}
      \textbf{(b)} Name every parent function whose domain is \emph{not} all real numbers.
      \writespace{1.0cm}{$y = \sqrt{x}$, $y = \frac{1}{x}$, $y = \frac{1}{x^2}$}
      \end{minipage}

\item Use even and odd degree. Which parent graph does $y = x^6$ look like, and which does
      $y = x^7$ look like? Describe where the arms of each go.
      \writespace{1.1cm}{$x^6$ is even degree, so it looks like $x^2$: both arms up. $x^7$ is odd, so it looks like $x^3$: left arm down, right arm up.}

\end{enumerate}

\vspace{0.06in}

\begin{headlinebox}{lilac}
\textbf{\color{plum}Part C --- Absolute value as a piecewise function}
\end{headlinebox}

\begin{enumerate}[leftmargin=1.9em, itemsep=7pt, start=8]

\item \textbf{(a)} Write $|x - 4|$ as a piecewise function. Where is the break, and why
      there?
      \writespace{1.4cm}{$|x - 4| = -(x - 4) = 4 - x$ for $x < 4$, and $x - 4$ for $x \ge 4$. The break is at $x = 4$, where $x - 4$ changes sign.}

      \textbf{(b)} Use your rules to find $|x - 4|$ at $x = 1$ and at $x = 7$. Name the rule
      you used each time.
      \writespace{1.2cm}{$1 < 4$, first rule: $4 - 1 = 3$. \quad $7 \ge 4$, second rule: $7 - 4 = 3$.}

\item Let $g(x) = |x + 1|$. Write $g$ as a piecewise function (where does $x + 1$ change
      sign?), then find $g(-4)$ using the rule that owns $x = -4$.
      \writespace{2.0cm}{$g(x) = -(x + 1)$ for $x < -1$, and $x + 1$ for $x \ge -1$. \par $-4 < -1$, so the first rule: $g(-4) = -(-4 + 1) = -(-3) = 3$.}

\end{enumerate}

\vspace{0.06in}

\begin{headlinebox}{lilac}
\textbf{\color{plum}Part D --- The headline, in writing}
\end{headlinebox}

\begin{enumerate}[leftmargin=1.9em, itemsep=7pt, start=10]

\item \textbf{(a)} A classmate says the rule $|x| = -x$ for $x < 0$ must be a typo, because
      ``$-x$ is negative.'' Use the words \emph{opposite} and \emph{input} to explain why the
      rule is right. Test one value to back it up.
      \writespace{1.6cm}{The rule $-x$ only receives inputs below 0, and $-x$ means the opposite of the input. The opposite of a negative is positive: at $x = -5$, $-x = -(-5) = 5 = |-5|$.}

      \textbf{(b)} The same classmate writes $\sqrt{-4}\cdot\sqrt{-9} = \sqrt{36} = 6$. Explain
      the mistake and give the right value.
      \writespace{1.2cm}{$\sqrt{a}\cdot\sqrt{b} = \sqrt{ab}$ needs $a, b \ge 0$. Rewrite with $i$ first: $(2i)(3i) = 6i^2 = -6$.}

\end{enumerate}

\boxguard[5]
\begin{spiralbox}
\small
\textbf{Next lesson (2.2, Transformations of Functions):} every graph in today's library can
be moved, stretched, and flipped. In problem 8, the corner of $|x|$ moved from $(0, 0)$ to
$(4, 0)$ --- next time we read moves like that straight from the equation.
\end{spiralbox}

\end{document}
